\documentclass{itmo}

\setlist[itemize]{label=---}
\setlist[enumerate]{label=\arabic*)}

\usepackage{dcolumn}
\newcolumntype{d}[1]{D{.}{.}{#1}}

\begin{document}

\makereporttitle{}
	{лабораторной работе №1}
	{Структуры и алгоритмы в базах данных и распределенных системах}
	{Реализация алгоритмов хэширования (perfect, extendible, min)}
	{}
	{Постнов~С.~А./P4135}
	{Платонов~А.~В./доцент, Портнов~П.~В./ассистент}

\maketableofcontents

\chapter{Теоретическая часть}

\section{Perfect hashing}

\textit{Perfect hash} для множества ключей $S$~---~это хэш--функция $h$, которая отображает ключи из $S$ в индексы таблицы без коллизий.
В классическом варианте (FKS) используется двухуровневая схема: первичная таблица распределяет ключи по корзинам, а для каждой корзины строится вторичная таблица, в которой подбираются параметры так, чтобы коллизий не было.
Основными идеями алгоритма являются:
\begin{itemize}
	\item первичная таблица имеет размер, равный ближайшему простому $\ge n$;
	\item в ячейке хранится либо одна пара \texttt{key/value}, либо указатель на вторичную таблицу;
	\item при коллизии создаётся/перестраивается вторичная таблица: увеличивается её размер и перебираются различные хэш--функции до получения раскладки без коллизий.
\end{itemize}

\section{Extendible hashing}

\textit{Extendible hashing} предназначен для динамических наборов данных, где число ключей заранее неизвестно и структура должна расти и уменьшаться без полной перестройки таблицы.
В отличие от \textit{perfect hashing}, здесь допускаются коллизии, но они локализуются в бакетах фиксированной ёмкости.
Основные элементы алгоритма:
\begin{itemize}
	\item директория содержит $2^d$ указателей на бакеты, где $d$ --- глобальная глубина;
	\item каждый бакет имеет локальную глубину, определяющую, по скольким младшим битам хэша он адресуется;
	\item при переполнении бакет расщепляется, а при необходимости директория удваивается;
	\item при удалении возможны обратные операции: слияние соседних бакетов и уменьшение директории.
\end{itemize}
Такая схема обеспечивает амортизированно быстрые операции \texttt{insert/get/update/delete} и хорошо подходит для внешней памяти, так как бакеты можно размещать непосредственно в файле и адресовывать через отображение файла в память.

\section{Min hash}

\textit{Min hash} предназначен для эффективного поиска похожих объектов в больших коллекциях, используя хэширование, чувствительное к локальности.
Мера сходства --- коэффициент Жаккара $J(A,B) = |A \cap B| / |A \cup B|$, где $A$ и $B$~---~множества \textit{шинглов} (подстрок длины $k$, извлечённых скользящим окном) двух документов.
\textit{Min hash} даёт компактную аппроксимацию Жаккара, то есть для набора хэш--функций $h_1,\ldots,h_m$ сигнатура документа --- это вектор минимумов $\sigma_i = \min_{s \in A} h_i(s)$.
Вероятность совпадения $\sigma_i(A) = \sigma_i(B)$ равна $J(A,B)$, поэтому долю совпавших компонент можно использовать как оценку сходства.

\textit{Locality--Sensitive Hashing} (LSH) ускоряет поиск кандидатов --- сигнатура делится на $b$ полос по $r$ элементов каждая.
Два документа попадают в один бакет хотя бы одной полосы с вероятностью $1 - (1 - s^r)^b$, где $s$~---~истинный Жаккар.
Порог чувствительности $s^* \approx (1/b)^{1/r}$ задаётся выбором $b$ и $r$.
Полосовая схема позволяет за линейное время построить список пар--кандидатов, а затем верифицировать их точным сравнением шинглов.

\chapter{Практическая часть}

\section{Реализация perfect hash таблицы}

Публичный интерфейс реализован в файлах \texttt{inc/perfect.h} и \texttt{src/perfect.c}:
\begin{itemize}
	\item \texttt{create\_perfect\_table} создаёт таблицу нужного размера;
	\item \texttt{perfect\_table\_from\_csv} строит таблицу по входному CSV;
	\item \texttt{perfect\_get} выполняет поиск значения по ключу;
	\item \texttt{free\_perfect\_table} освобождает память.
\end{itemize}

Структура \texttt{perfect\_hash\_table\_t} содержит массив указателей на \texttt{ceil\_value\_t}.
Каждая ячейка либо хранит одну пару \texttt{key/value}, либо указывает на вторичную таблицу (вложенный \texttt{perfect\_hash\_table\_t}).
Для хэширования строк используются функции из \texttt{inc/hash.h}.
При перестроении вторичных таблиц перебирается набор хэш--функций до построения раскладки без коллизий.

\section{Реализация extendible hash таблицы}

Основной код расположен в файле \texttt{extendible/extendible.go}.
Таблица хранится не в оперативной памяти как отдельная структура с последующей сериализацией, а непосредственно в файле, отображённом в память с помощью \texttt{mmap}.
Файл разбит на три логические области:
\begin{enumerate}
	\item заголовок с метаданными таблицы (магическое число, версия формата, глобальная глубина, размер бакета);
	\item директория фиксированного максимального размера, содержащая идентификаторы бакетов;
	\item область бакетов, где для каждого бакета хранятся локальная глубина, число записей и массив пар \texttt{key/value/hash}.
\end{enumerate}

Публичный интерфейс включает операции \texttt{Insert}, \texttt{Update}, \texttt{Delete}, \texttt{Get}, а также открытие и закрытие таблицы в файле.
При переполнении вызывается split бакета, при удалении --- merge buddy-бакетов и, если возможно, shrink директории.

\section{Реализация min hash таблицы LSH индекса}

Основные файлы: \texttt{min/minhash.go} (ядро алгоритма) и \texttt{min/benchmark\_test.go} (генерация случайного корпуса для замеров).
Архитектура индекса:
\begin{enumerate}
	\item из текста документа извлекаются шинглы функцией \texttt{shingleHashes} скользящим окном длины $k=2$ слова; значение каждого шингла --- FNV-1a хэш строки;
	\item функцией \texttt{signatureFor} строится min hash--сигнатура из $m=64$ хэш--функций на основе splitmix64;
	\item при построении индекса (\texttt{Build}) сигнатуры всех документов разбиваются на $b=8$ полос по $r=8$ компонент; каждая полоса хэшируется FNV--полиномом в ключ бакета (\texttt{bandKey});
	\item при добавлении нового документа (\texttt{Add}) его бакеты объединяются с существующими;
	\item поиск дубликатов (\texttt{FindDuplicates}) извлекает кандидатов из бакетов и проверяет пары точным подсчётом Жаккара на сохранённых шинглах;
	\item \texttt{FullScanDuplicates} является базовым алгоритмом сравнения: каждый документ проверяется против всех остальных за $O(n^2)$.
\end{enumerate}
Параметры $m$, $b$, $r$, $k$ передаются через структуру \texttt{Config}.

\chapter{Исследовательская часть}

\section{Аппаратные характеристики}

Все замеры проводились на ноутбуке со следующей конфигурацией:
\begin{itemize}
	\item операционная система: \texttt{macOS 26.3.1} (arm64);
	\item процессор: \texttt{Apple M3 Pro}, 11 ядер;
	\item оперативная память: \texttt{18 ГБ};
	\item версия Go (для \texttt{extendible} и \texttt{min}): \texttt{go1.25}, версия C (для \texttt{perfect}): \texttt{с17};
\end{itemize}
Во время замеров не выполнялись ресурсоёмкие фоновые задачи, а ноутбук был подключён к сети электропитания для обеспечения максимальной производительности.

\section{Методика замеров perfect hashing}

Для замеров использована программа \texttt{src/benchmark.c}.
Для каждого датасета из \texttt{data/bench/manifest.txt}:
\begin{enumerate}
	\item создаётся perfect hash таблица из csv датасета;
	\item измеряется время чтения всех $n$ ключей (\texttt{perfect\_get}).
\end{enumerate}
Каждый прогон усредняется по \texttt{BENCH\_ITERS} (по умолчанию 100).

\section{Методика замеров extendible hashing}

Для extendible hashing использованы встроенные benchmark-тесты Go и pprof.
Измерялись четыре операции: \texttt{insert}, \texttt{update}, \texttt{delete}, \texttt{get}.
Для каждой мощности набора данных выполнялись бенчмарки со считыванием времени на операцию, пропускной способности, количества аллокаций и объёма выделенной памяти.

\section{Методика замеров min hash}

Для min hash использованы встроенные benchmark-тесты Go.
Измерялись три операции:
\begin{itemize}
	\item \texttt{build}~---~построение индекса из корпуса размером $N$;
	\item \texttt{add}~---~добавление $N$ документов по одному в уже существующий индекс;
	\item \texttt{fullscan}~---~полный обход: для каждого из $N$ документов запускается линейный поиск дубликатов по всему корпусу.
\end{itemize}

Для каждой мощности набора данных фиксировались время на элемент, пропускная способность, количество итераций и 95\% доверительный интервал.

\section{Результаты для perfect hashing}

Таблица~\ref{tab:perfect-bench} содержит усреднённые значения задержки поиска и пропускной способности.
\begin{table}[h]
	\centering
	\caption{Результаты бенчмарка perfect hash (среднее по итерациям)}
	\label{tab:perfect-bench}
	\begin{tabularx}{\textwidth}{|X|X|X|}
		\toprule
		$N$ & поиск, нс/оп & поиск, млн оп/с \\
		\midrule
		20000 & $23.1 \pm 0.7$ & 44.19 \\
		50000 & $22.1 \pm 0.1$ & 45.33 \\
		100000 & $22.2 \pm 0.1$ & 45.07 \\
		250000 & $24.6 \pm 0.5$ & 40.94 \\
		500000 & $28.0 \pm 0.7$ & 36.08 \\
		750000 & $26.7 \pm 0.7$ & 37.77 \\
		1100000 & $29.9 \pm 0.7$ & 33.63 \\
		2000000 & $30.5 \pm 0.7$ & 33.06 \\
		5000000 & $34.3 \pm 0.9$ & 29.36 \\
		7500000 & $36.1 \pm 1.0$ & 27.94 \\
		\bottomrule
	\end{tabularx}
\end{table}

Графики результатов замеров представлены на рисунках~\ref{img:perfect_bench_avg_latency} и~\ref{img:perfect_bench_throughput}:
\includeimage
	{perfect_bench_avg_latency}
	{f}
	{H}
	{\textwidth}
	{Средняя задержка операции поиска для perfect hash}

\includeimage
	{perfect_bench_throughput}
	{f}
	{H}
	{\textwidth}
	{Пропускная способность операции поиска для perfect hash}

По таблице~\ref{tab:perfect-bench} и графикам можно сделать вывод, что при росте $N$ задержка \texttt{get} слабо растет, примерно с 22--23 нс/оп на малых наборах до 34--36 нс/оп на максимальных.
Пропускная способность, соответственно, снижается с порядка 45 до 28 млн оп/с, что ожидаемо из-за роста рабочей области и нагрузки на кэш.
В целом для статического множества ключей perfect hashing сохраняет высокую производительность поиска и предсказуемое время доступа.

\section{Результаты для extendible hashing}

Таблицы~\ref{tab:extendible-latency}, \ref{tab:extendible-throughput} и~\ref{tab:extendible-warmup} содержат результаты бенчмарка extendible hashing для всех измеренных размеров набора данных.
\begin{table}[H]
	\centering
	\caption{Результаты бенчмарка задержки в extendible hashing (среднее по итерациям)}
	\label{tab:extendible-latency}
	\begin{tabularx}{\textwidth}{|l|l|X|X|X|X|}
		\toprule
		$N$ & итерации & вставка, нс/оп & обновление, нс/оп & удаление, нс/оп & поиск, нс/оп \\
		\midrule
		1000 & 27894 & $43.42 \pm 0.035$ & $32.61 \pm 0.019$ & $40.86 \pm 0.024$ & $32.41 \pm 0.020$ \\
		5000 & 5374 & $44.61 \pm 0.060$ & $25.88 \pm 0.023$ & $35.37 \pm 0.032$ & $25.16 \pm 0.023$ \\
		10000 & 2610 & $45.85 \pm 0.072$ & $27.47 \pm 0.028$ & $40.36 \pm 0.048$ & $26.73 \pm 0.028$ \\
		50000 & 483 & $49.59 \pm 0.109$ & $30.86 \pm 0.053$ & $70.55 \pm 0.355$ & $30.81 \pm 0.046$ \\
		100000 & 208 & $57.37 \pm 0.133$ & $32.03 \pm 0.063$ & $94.91 \pm 0.861$ & $31.87 \pm 0.051$ \\
		500000 & 19 & $123.7 \pm 0.699$ & $36.31 \pm 0.103$ & $310.8 \pm 6.781$ & $35.95 \pm 0.087$ \\
		1000000 & 10 & $214.1 \pm 2.413$ & $54.02 \pm 0.440$ & $559.7 \pm 19.71$ & $53.44 \pm 0.416$ \\
		\bottomrule
	\end{tabularx}
\end{table}

\begin{table}[H]
	\centering
	\caption{Результаты бенчмарка пропускной способности в extendible hashing (среднее по итерациям)}
	\label{tab:extendible-throughput}
	\begin{tabularx}{\textwidth}{|X|X|X|X|X|X|}
		\toprule
		$N$ & итерации & вставка, млн оп/с & обновление, млн оп/с & удаление, млн оп/с & поиск, млн оп/с \\
		\midrule
		1000 & 27894 & $23.0 \pm 0.018$ & $30.7 \pm 0.018$ & $24.5 \pm 0.015$ & $30.9 \pm 0.019$ \\
		5000 & 5374 & $22.4 \pm 0.030$ & $38.6 \pm 0.021$ & $28.3 \pm 0.025$ & $39.7 \pm 0.036$ \\
		10000 & 2610 & $21.8 \pm 0.034$ & $36.4 \pm 0.029$ & $24.8 \pm 0.047$ & $37.4 \pm 0.035$ \\
		50000 & 483 & $20.2 \pm 0.044$ & $32.4 \pm 0.053$ & $14.2 \pm 0.071$ & $32.5 \pm 0.047$ \\
		100000 & 208 & $17.4 \pm 0.051$ & $31.2 \pm 0.061$ & $10.5 \pm 0.097$ & $31.4 \pm 0.051$ \\
		500000 & 19 & $8.1 \pm 0.069$ & $27.5 \pm 0.098$ & $3.2 \pm 0.035$ & $27.8 \pm 0.083$ \\
		1000000 & 10 & $4.7 \pm 0.048$ & $18.5 \pm 0.151$ & $1.8 \pm 0.063$ & $18.7 \pm 0.146$ \\
		\bottomrule
	\end{tabularx}
\end{table}

\begin{table}[H]
	\centering
	\caption{Результаты замера времени прогрева (warmup) в extendible hashing}
	\label{tab:extendible-warmup}
	\begin{tabularx}{\textwidth}{|X|X|X|}
		\toprule
		$N$ & итерации & время warmup, нс/оп \\
		\midrule
		1000 & 1681 & $815416 \pm 128.5$ \\
		5000 & 8356 & $1093872 \pm 72.15$ \\
		10000 & 1124 & $1090091 \pm 196.1$ \\
		50000 & 1166 & $1062745 \pm 170.0$ \\
		100000 & 283 & $3677399 \pm 588.4$ \\
		500000 & 250 & $6045153 \pm 653.3$ \\
		1000000 & 100 & $10683703 \pm 1008$ \\
		\bottomrule
	\end{tabularx}
\end{table}

Из таблиц можно сделать вывод, что операции \texttt{get} и \texttt{update} наиболее стабильны по мере роста набора данных.
Операция \texttt{insert} заметно дороже из-за split бакетов и расширения директории, а \texttt{delete} оказывается самой тяжёлой на больших $N$ из-за merge/shrink после удаления элементов.
Дополнительные профилировочные замеры показывают, что основной объём аллокаций приходится на \texttt{insert}; для \texttt{update}, \texttt{delete} и \texttt{get} объём дополнительных аллокаций почти не растёт и определяется в основном накладными расходами рантайма.

Графики результатов замеров extendible hashing представлены на рисунках ниже.
\includeimage
	{extendible_benchmark_latency}
	{f}
	{H}
	{\textwidth}
	{Задержка операций extendible hashing на элемент}

\includeimage
	{extendible_benchmark_throughput}
	{f}
	{H}
	{\textwidth}
	{Пропускная способность операций extendible hashing}

\includeimage
	{extendible_warmup_elapsed_ns}
	{f}
	{H}
	{\textwidth}
	{Время прогрева (warmup) для extendible hashing}

\section{Результаты для min hash}

Таблица~\ref{tab:min-latency} содержит задержку операций, таблица~\ref{tab:min-throughput}~---~пропускную способность.
\begin{table}[H]
	\centering
	\caption{Результаты бенчмарка задержки в min hash (среднее по итерациям)}
	\label{tab:min-latency}
	\begin{tabularx}{\textwidth}{|l|X|X|X|}
		\toprule
		$N$ & построение, тыс нс/элемент & добавление, тыс нс/элемент & полный обход, тыс нс/элемент \\
		\midrule
		500 & $14.940 \pm 0.008846$ ($n=9280$) & $15.329 \pm 0.004471$ ($n=9378$) & $182.112 \pm 0.1235$ ($n=789$) \\
		1000 & $15.164 \pm 0.006001$ ($n=4753$) & $15.238 \pm 0.006451$ ($n=4718$) & $364.406 \pm 0.3443$ ($n=196$) \\
		2500 & $15.149 \pm 0.009802$ ($n=1902$) & $15.258 \pm 0.01479$ ($n=1891$) & $911.679 \pm 1.548$ ($n=31$) \\
		5000 & $15.145 \pm 0.01144$ ($n=950$) & $15.338 \pm 0.01724$ ($n=938$) & $1825.010 \pm 2.314$ ($n=7$) \\
		7500 & $15.242 \pm 0.01519$ ($n=628$) & $15.453 \pm 0.01588$ ($n=621$) & $2750.945 \pm 3.632$ ($n=3$) \\
		\bottomrule
	\end{tabularx}
\end{table}

\begin{table}[H]
	\centering
	\caption{Результаты бенчмарка пропускной способности в min hash (среднее по итерациям)}
	\label{tab:min-throughput}
	\begin{tabularx}{\textwidth}{|l|X|X|X|}
		\toprule
		$N$ & построение, оп/с & добавление, оп/с & полный обход, оп/с \\
		\midrule
		500 & $66936 \pm 40$ ($n=9280$) & $65234 \pm 19$ ($n=9378$) & $5491 \pm 4$ ($n=789$) \\
		1000 & $65944 \pm 26$ ($n=4753$) & $65624 \pm 28$ ($n=4718$) & $2744 \pm 3$ ($n=196$) \\
		2500 & $66012 \pm 43$ ($n=1902$) & $65538 \pm 64$ ($n=1891$) & $1097 \pm 2$ ($n=31$) \\
		5000 & $66029 \pm 50$ ($n=950$) & $65196 \pm 73$ ($n=938$) & $548 \pm 1$ ($n=7$) \\
		7500 & $65610 \pm 65$ ($n=628$) & $64711 \pm 66$ ($n=621$) & $364 \pm 1$ ($n=3$) \\
		\bottomrule
	\end{tabularx}
\end{table}

Операции \texttt{build} и \texttt{add} демонстрируют практически постоянную задержку (около $15\,\text{тыс нс/элемент}$) при росте корпуса, что объясняется независимой обработкой каждого документа при вычислении сигнатуры.
Задержка \texttt{fullscan} растёт квадратично ($O(n^2)$), поскольку для каждого из $N$ документов выполняется линейный перебор остальных документов корпуса.

\chapter{Вывод}

Реализована perfect hash таблица для строковых ключей в статическом сценарии (построение по CSV и поиск), extendible hash таблица с файловым хранением на базе \texttt{mmap}, поддерживающая динамические операции вставки, чтения, обновления и удаления, а также min hash LSH индекс для поиска почти дубликатов в текстовых коллекциях.

Проведены замеры производительности на наборах данных разного размера.
Для perfect hashing подтверждена эффективность на статическом множестве ключей: операция \texttt{get} показывает низкую задержку и высокую пропускную способность на всем диапазоне размеров данных.
Для extendible hashing показано, что операции \texttt{get} и \texttt{update} остаются самыми быстрыми, тогда как \texttt{insert} и особенно \texttt{delete} несут дополнительные издержки на модификацию структуры каталогов и бакетов.
Для min hash подтверждена линейная масштабируемость операций \texttt{build} и \texttt{add} ($\approx 15\,\text{тыс нс/элемент}$) при квадратичном росте задержки \texttt{fullscan}~---~базового алгоритма перебора, используемого для эталонного сравнения.
Результаты представлены в таблицах и на графиках.

\end{document}
